\documentclass[../../../uem_phy_std_q.tex]{subfiles}

\begin{document}
    \begin{enumerate}
        \item%
            図のように，電池E（起電力$V_0$），電気抵抗$\mathrm{R}_1$（抵抗値$2R$），%
            $\mathrm{R}_2$（抵抗値$3R$）からなる回路がある．AB間の合成抵抗$R_\mathrm{T}$を求めよ．\\
            \begin{minipage}{\linewidth}
                \centering
                    \includegraphics{\subfix{fig_uem_phy_std_em_cx_03/fig_uem_phy_std_em_cx_03_01_q.pdf}}
            \end{minipage}
            \vspace{1pt}
        \item%
            図のように，電池E（起電力$V_0$），%
            電気抵抗$\mathrm{R}_1$（抵抗値$2R$），$\mathrm{R}_2$（抵抗値$3R$）からなる回路がある．%
            AB間の合成抵抗$R_\mathrm{T}$を求めよ．\\
            \begin{minipage}{\linewidth}
                \centering
                    \includegraphics{\subfix{fig_uem_phy_std_em_cx_03/fig_uem_phy_std_em_cx_03_02_q.pdf}}
            \end{minipage}
            \vspace{1pt}
        \item%
            図のように，電池E$_1$（起電力$2V_0$），E$_2$（起電力$V_0$），%
            電気抵抗$\mathrm{R}_1$（抵抗値$2R$），$\mathrm{R}_2$（抵抗値$R$），%
            $\mathrm{R}_3$（抵抗値$2R$）からなる回路がある．%
            AB間の合成抵抗$R_\mathrm{T}$を求めよ．\\
            \begin{minipage}{\linewidth}
                \centering
                    \includegraphics{\subfix{fig_uem_phy_std_em_cx_03/fig_uem_phy_std_em_cx_03_03_q.pdf}}
            \end{minipage}
            \vspace{1pt}
        \needspace{8\baselineskip}
        \item%
            図のように，電気抵抗$\mathrm{R}_1$，$\mathrm{R}_2$，$\mathrm{R}_3$，%
            $\mathrm{R}_4$，$\mathrm{R}_5$からなる回路がある．%
            $\mathrm{R}_1$，$\mathrm{R}_4$，$\mathrm{R}_5$の抵抗値は$R$であり，%
            $\mathrm{R}_2$，$\mathrm{R}_3$の抵抗値は$2R$である．AB間の合成抵抗$R_\mathrm{T}$を求めよ．\\
            \begin{minipage}{\linewidth}
                \centering
                    \includegraphics{\subfix{fig_uem_phy_std_em_cx_03/fig_uem_phy_std_em_cx_03_04_q.pdf}}
            \end{minipage}
            \vspace{1pt}
        \needspace{10\baselineskip}
        \item%
            図のように，12個の電気抵抗からなる回路がある．いずれの電気抵抗の抵抗値も$R$である．%
            AB間の合成抵抗$R_\mathrm{AB}$を求めよ．\\
            \begin{minipage}{\linewidth}
                \centering
                    \includegraphics{\subfix{fig_uem_phy_std_em_cx_03/fig_uem_phy_std_em_cx_03_05_q.pdf}}
            \end{minipage}
            \vspace{1pt}
    \end{enumerate}
\end{document}




